\documentclass[10pt, a4paper]{article}

% --- PACKAGES ---
\usepackage[utf8]{inputenc}
\usepackage[T1]{fontenc}
\usepackage[french]{babel}
\usepackage[defaultsans]{lato}
\renewcommand{\familydefault}{\sfdefault}
\usepackage[a4paper, margin=0.55in]{geometry}
\usepackage{enumitem}
\setlist[itemize]{leftmargin=*, nosep, itemsep=2pt, after=\vspace{0.3\baselineskip}}
\usepackage{titlesec}
\usepackage{parskip}
\usepackage{xcolor}
\usepackage{fontawesome5}
\usepackage{hyperref}
\usepackage{array}
\usepackage{graphicx}

% --- CONFIGURATION ---
\newcolumntype{L}{>{\raggedright\arraybackslash}p{0.48\textwidth}}
\newcolumntype{K}{>{\raggedright\arraybackslash}p{0.40\textwidth}}
\newcolumntype{P}{>{\raggedright\arraybackslash}p{0.25\textwidth}}
\newcolumntype{G}{>{\raggedright\arraybackslash}p{0.73\textwidth}}

\definecolor{accent}{RGB}{20, 60, 110}
\hypersetup{colorlinks=true, urlcolor=accent, linkcolor=accent}

\titleformat{\section}{\large\bfseries\color{accent}}{}{0em}{}[\titlerule]
\titlespacing{\section}{0pt}{12pt}{6pt}

\begin{document}
	
	% --- HEADER ---
	\noindent
	\begin{minipage}[t]{0.20\textwidth}
		\vspace{0pt}
		% Assure-toi que "photoprofil.jpg" est bien dans le dossier
		\includegraphics[width=\textwidth]{photoprofil.jpg}
	\end{minipage}
	\hfill
	\begin{minipage}[t]{0.75\textwidth}
		\vspace{0pt}
		{\Huge \bfseries \color{accent} AURÉLIEN LOGEAIS} \\
		{\large \textbf{Analyste Cybersécurité Junior | Alternant}} \\
		\textbf{\color{accent}Disponible pour un CDI dès Septembre 2026} \\
		
		\vspace{8pt}
		
		\begin{tabular}{@{}K K@{}}
			\faEnvelope \hspace{2pt} \href{mailto:aurelienlogeais781@gmail.com}{aurelienlogeais781@gmail.com} & \faPhone \hspace{2pt} 07 87 32 05 72 \\[1ex]
			\faLinkedin \hspace{2pt} \href{https://www.linkedin.com/in/aurelien-logeais-1727922b4}{in/aurelienlogeais} & \faMapMarker* \hspace{2pt} Marseille / Aix-en-Provence \\[1ex]
			\faGlobe \hspace{2pt} \href{https://portfolio.pmvix.com}{CV Interactif} & \faCar \hspace{2pt} Permis B
		\end{tabular}
	\end{minipage}
	
	\vspace{10pt}
	
	% --- PROFIL ---
	\section*{Profil}
	Actuellement en dernière année de Mastère Cybersécurité à l'ESGI. Je combine une solide base théorique validée par la certification Security+ avec une expertise pratique acquise en alternance sur Azure et IAM, ainsi que via un laboratoire de sécurité personnel avancé centré sur le SIEM et l'automatisation. Passionné par la Blue Team, je recherche un poste d'Analyste SOC ou d'Ingénieur Cybersécurité orienté détection et automatisation.
	
	% --- CERTIFICATIONS & LAB ---
	\section*{Certifications \& Laboratoire (Homelab)}
	
	\noindent
	\begin{tabular}{@{}ll@{}}
		\textbf{CompTIA Security+ :} & \textbf{Validée depuis Janvier 2026} \\
		\textbf{Stormshield CSNA :} & \textbf{En cours de préparation pour un passage en Juillet 2026}
	\end{tabular}
	
	\vspace{6pt}
	\noindent \textbf{Infrastructure Homelab (Proxmox) :} Conception d'une architecture sécurisée complète.
	\vspace{-4pt}
	\begin{itemize}[label={--}, leftmargin=1.0cm, topsep=0pt, partopsep=0pt]
		\item \textbf{Détection et SIEM :} Déploiement complet de Splunk Enterprise avec création de règles de détection personnalisées et ingestion de logs.
		\item \textbf{Automatisation :} Écriture de playbooks Ansible et développement de scripts Bash ou Python pour assurer le déploiement.
		\item \textbf{Sécurité Réseau :} Segmentation stricte via pfSense utilisant des VLANs et filtrage DNS via AdGuard.
		\item \textbf{Accès Distant :} Mise en place d'une architecture Zero Trust via Cloudflare Tunnel et VPN Wireguard.
	\end{itemize}
	
	% --- EXPÉRIENCES ---
	\section*{Expériences Professionnelles}
	\textbf{Assistant Business App \& Sécurité Cloud (Alternance)} \hfill \textbf{Août 2024 -- Septembre 2026} \\
	\textit{Kaizzen, Marseille}
	\begin{itemize}
		\item Gestion des identités et des accès IAM sur l'environnement Microsoft Azure et Entra ID.
		\item Sécurisation des secrets applicatifs et automatisation des processus métier.
		\item Veille technologique et application des bonnes pratiques de sécurité Cloud.
	\end{itemize}
	
	\textbf{Responsable d'équipe} \hfill \textbf{Étés 2022 -- 2024} \\
	\textit{Plage \& Loisirs, Le Lavandou}
	\begin{itemize}
		\item Management direct et coordination du personnel dans un environnement à forte affluence.
		\item Prise de décision rapide, résolution de problèmes en temps réel et maintien des standards de qualité.
	\end{itemize}
	
	% --- PROJETS ACADÉMIQUES ---
	\section*{Projets Académiques}
	\begin{tabular}{@{}P G@{}}
		\textbf{Phishing Tool \& IA :} & Développement d'une plateforme de simulation d'attaques assistée par IA.\\
		\textbf{Blockchain :} & Création d'un protocole pour la certification immuable de diplômes. \\
		\textbf{Architecture :} & Développement bas niveau incluant un Moteur Ray Tracing et un Compilateur RAM.
	\end{tabular}
	
	% --- COMPÉTENCES ---
	\section*{Compétences Techniques}
	\begin{tabular}{@{}P G@{}}
		\textbf{SOC \& Détection :} & Splunk, création de règles d'alerte, analyse de logs et framework MITRE ATT\&CK. \\
		\textbf{Automatisation :} & Développement en Bash, PowerShell et Python et playbooks Ansible. \\
		\textbf{Réseau \& Défense :} & Administration de pfSense, analyse Wireshark, scan Nmap et Cloudflare Zero Trust. \\
		\textbf{Systèmes \& Cloud :} & Administration Linux Debian, Windows Server, conteneurisation Docker et Azure AD.
	\end{tabular}
	
	% --- FORMATION ---
	\section*{Formation}
	\noindent
	\begin{tabular*}{\textwidth}{@{}l@{\extracolsep{\fill}}r@{}}
		\textbf{Mastère Cybersécurité} | ESGI Aix-en-Provence & \textbf{2024 -- 2026} \\
		\textbf{Licence Informatique (Label Cyberedu)} | Université de Toulon & \textbf{2021 -- 2024} \\
	\end{tabular*}
	
	
\end{document}